\begin{Huge}
    Medizininformatik-FAQ
\end{Huge}

\begin{block}{H\"aufig gestellte Fragen zum Studium}

\question{Lernt man im Studium, wie man programmiert?}{
    Ja, aber auf einer sehr eigenst\"andigen Basis. Man bekommt einen \"Uberblick \"uber die Sprache(n) und \"uber dahinterstehende Programmierkonzepte. Praktische Programmiererfahrung muss man sich aber gr\"o"stenteils selbst erarbeiten.
}
\question{Welche Programmiersprachen macht man da so?}{
    Das ist von Prof zu Prof und von Veranstaltung zu Veranstaltung unterschiedlich. In den ersten beiden Semestern lernt man meistens Racket und Java.
}
\question{Muss man programmieren k\"onnen, um das Studium anzufangen?}{
    Nein. Die Vorlesung beginnt absolut bei 0, um allen den Einstieg zu erm\"oglichen.
}

\question{Muss man gut in Mathe sein?}{
    Man muss kein Mathe-Genie sein, man sollte Mathe aber nicht hassen. Es ist gerade am Anfang viel Mathe.
}
\question{Ich will eigentlich Medizin studieren, aber mein NC reicht nicht. Medizininformatik ist doch auch was mit Medizin, oder?}{
    Nein, Medizininformatik ersetzt kein Medizinstudium. Man lernt zwar Grundlagen der Anatomie, Histologie und Pathologie, aber auf einem grundlegenderen Niveau wie in einem Medizinstudium. Deshalb h\"ort man auch andere Vorlesungen wie die normalen Medizin-Studis. Zweck dieser Vorlesungen ist, am Ende ungef\"ahr zu verstehen, worum es geht, wenn man Anwendungen f\"ur Mediziner*innen baut.
}

\question{Wie ist die Frauenquote so?}{
    72\%.
}
\question{Was ist der Unterschied zwischen Bio- und Medizininformatik?}{
    Die Bioinformatik besch\"aftigt sich grob gesagt mit automatisierter Verarbeitung von DNA, Molek\"ulstrukturen etc., Medizininformatik geht mehr in Richtung Patientendaten und medizinische Bildverarbeitung.
}
\question{Gibt es Praktika?}{
    Im normalen Studienverlauf ist kein berufsorientiertes Praktikum vorgesehen, viele arbeiten aber parallel als Werksstudent*in oder man macht ein Kurzpraktikum in den Semesterferien.
}

\question{Kann man ein Auslandssemester machen?}{
    Klar, geht immer. T\"ubingen nimmt am ERASMUS-Programm teil, die Organisation ist aber langwierig und man sollte sich fr\"uh drum k\"ummern.
}

\question{Was arbeitet man danach so?}{
    Alle Bereiche der IT-Branche, insbesondere in den vielf\"altigen Berufsfeldern der medizinischen Informationsverarbeitung und des Gesundheitswesens.
}

\question{Wie ist da so der NC?}{
    Die letzten 4 Jahre gab es keinen NC. Das muss aber theoretisch nicht immer so sein.
}
\end{block}
